% Part: first-order-logic
% Chapter: syntax-and-semantics
% Section: satisfaction

\documentclass[../../../include/open-logic-section]{subfiles}

\begin{document}

\olfileid{fol}{syn}{sat}

\olsection{\article{structure} \printtoken{S}{structure} मधील
  \article{formula} \printtoken{S}{formula} ची पूर्ति}

\begin{explain}
एकीकडे पदे आणि !!{formula}s यांसारख्या अभिव्यक्ती आणि दुसरीकडे
!!{structure}s यांना जोडणाऱ्या मूलभूत संकल्पना म्हणजे पदाचे
\emph{!!{value}} आणि !!a{formula} ची \emph{पूर्ति}. अनौपचारिकपणे,
पदाचे !!{value} म्हणजे !!a{structure} मधील !!{element} असतो---पद
नुसते स्थिरचिन्ह असेल, तर त्या स्थिरचिन्हाला !!{structure} ने नेमून
दिलेली वस्तू हे त्याचे !!{value} असते; आणि ते !!{function}s वापरून
बनवलेले असेल, तर पदातील स्थिरचिन्हांची !!{value}s आणि फलनचिन्हांना
नेमून दिलेली फलने यांपासून त्याचे !!{value} काढले जाते. विधेयांना
दिलेल्या अर्थनिर्धारणामुळे !!{structure} च्या प्रांतात !!{formula} सत्य
होत असेल, तर !!^a{formula} ची त्या !!a{structure} मध्ये \emph{पूर्ति
होते}. पूर्तिची ही संकल्पना विगमनाधारित रीतीने विनिर्दिष्ट केली जाते:
आण्विक !!{formula}s ची पूर्ति केव्हा होते हे !!{structure} चे
विनिर्देशन थेट सांगते; आणि एखाद्या जटिल !!{formula} ची पूर्ति केव्हा
होते हे आपण तिचा मुख्य संयोजक किंवा संख्यापक आणि तिच्या तात्काळ
!!{subformula}s ची पूर्ति होते की नाही यांवरून परिभाषित करतो.

इथे संख्यापकांची स्थिती थोडी अवघड आहे, कारण संख्यकीकृत !!{formula}
च्या तात्काळ !!{subformula} मध्ये एक मुक्त !!{variable} असते आणि
!!{structure}s ही !!{variable}s ची !!{value}s विनिर्दिष्ट करत नाहीत.
ही अडचण हाताळण्यासाठी आपण \emph{चर-मूल्यांकने} ही आणतो आणि पूर्तिची
व्याख्या केवळ !!a{structure} च्या सापेक्ष न देता, !!a{structure} आणि
!!a{variable} चे मूल्यांकन या दोहोंच्या सापेक्ष देतो.
\end{explain}

\begin{defn}[चर-मूल्यांकन]
!!a{structure}~$\Struct{M}$ हिच्यासाठीचे \emph{चर-मूल्यांकन}~$s$
म्हणजे प्रत्येक !!{variable} चे~$\Domain M$ मधील !!a{element} कडे
प्रतिचित्रण करणारे फलन; म्हणजेच, $s\colon \Var \to \Domain M$.
\end{defn}

\begin{explain}
!!^a{structure} प्रत्येक !!{constant} ला !!a{value} नेमून देते, तर
चर-मूल्यांकन प्रत्येक चराला मूल्य नेमून देते. पण त्यांच्यापासून बनलेली
पदेही !!{domain} मधील !!{element}s ना नावे देण्यासाठी वापरायची आहेत.
त्यासाठी आपण पदांचे !!{value} विगमनाधारित रीतीने परिभाषित करतो.
!!{constant}s आणि चरांसाठी मूल्य अनुक्रमे !!{structure} किंवा
चर-मूल्यांकन विनिर्दिष्ट करते तेच असते; अधिक जटिल पदांसाठी ते
!!{structure} ने !!{function}s ना नेमून दिलेली फलने वापरून पुनरावर्ती
रीतीने काढले जाते.
\end{explain}

\begin{defn}[पदांचे !!^{value}]
$t$ हे भाषा~$\Lang L$ हिचे पद, $\Struct M$ ही~$\Lang L$ साठीची
!!{structure} आणि $s$ हे~$\Struct M$ साठीचे !!a{variable} चे
मूल्यांकन असेल, तर \emph{!!{value}}~$\Value{t}{M}[s]$ पुढीलप्रमाणे
परिभाषित केले जाते:
\begin{enumerate}
\item \indcase{t}{c}{$\Value{\indfrm}{M}[s] = \Assign{\indcomplex}{M}$.}
\item \indcase{t}{x}{$\Value{\indfrm}{M}[s] = s(\indcomplex)$.}
\item \indcase{t}{\Atom{f}{t_1, \ldots, t_n}}{
\[
\Value{\indfrm}{M}[s] = \Assign{f}{M}(\Value{t_1}{M}[s], \ldots,
\Value{t_n}{M}[s]).
\]}
\end{enumerate}
\end{defn}

\begin{defn}[$x$-पर्याय]
$s$ हे !!a{structure}~$\Struct M$ साठीचे !!a{variable} चे मूल्यांकन
असेल, तर~$s$ पासून फार तर~$x$ ला नेमून दिलेल्या मूल्यातच भिन्न असणारे
$\Struct M$ साठीचे कोणतेही !!{variable} चे मूल्यांकन~$s'$ याला~$s$ चा
\emph{$x$-पर्याय} म्हणतात. $s'$ हा~$s$ चा $x$-पर्याय असेल, तर आपण
$\varAssign{s'}{s}{x}$ असे लिहितो.
\end{defn}

\begin{explain}
मूल्यांकन~$s$ च्या $x$-पर्यायाने~$x$ ला काही वेगळेच नेमून दिले
\emph{पाहिजे} असे नाही, हे लक्षात घ्या. वस्तुतः, प्रत्येक मूल्यांकन
स्वतःचाच $x$-पर्याय मानला जातो.
\end{explain}

\begin{defn}
  $s$ हे !!a{structure}~$\Struct M$ साठीचे !!a{variable} चे मूल्यांकन
  आणि $m \in \Domain{M}$ असेल, तर मूल्यांकन~$\Subst{s}{m}{x}$ हे पुढील
  रीतीने परिभाषित केलेले चर-मूल्यांकन आहे:
  \[\Subst{s}{m}{x}(y) = \begin{cases}
    m & \text{जर } y \ident x\\
    s(y) & \text{अन्यथा}.
  \end{cases}\]
\end{defn}

दुसऱ्या शब्दांत, $\Subst{s}{m}{x}$ हा~$s$ चा असा विशिष्ट $x$-पर्याय
आहे, जो प्रांतातील !!{element}~$m$ हे~$x$ ला नेमून देतो आणि~$x$ शिवाय
इतर !!{variable}s ना~$s$ ने नेमून दिलेल्याच गोष्टी नेमून देतो.

\begin{defn}[पूर्ति]
\ollabel{defn:satisfaction}
!!a{structure}~$\Struct M$ मध्ये !!a{variable} चे मूल्यांकन~$s$ याच्या
सापेक्ष !!a{formula}~$!A$ ची पूर्ति, संकेतांत $\Sat{M}{!A}[s]$, पुढीलप्रमाणे
पुनरावर्ती रीतीने परिभाषित केली जाते. (``$\Sat{M}{!A}[s]$ नाही'' या
अर्थाने आपण $\Sat/{M}{!A}[s]$ असे लिहितो.)
\begin{enumerate}
\tagitem{prvFalse}{%
  \indcase{!A}{\lfalse}{$\Sat/{M}{\indfrm}[s]$.}}{}

\tagitem{prvTrue}{%
  \indcase{!A}{\ltrue}{$\Sat{M}{\indfrm}[s]$.}}{}

\item \indcase{!A}{\Atom{R}{t_1, \dots, t_n}}{$\Sat{M}{\indfrm}[s]$
  तेव्हा आणि केवळ तेव्हाच, जेव्हा $\langle \Value{t_1}{M}[s], \dots,
  \Value{t_n}{M}[s] \rangle \in \Assign{R}{M}$.}

\item \indcase{!A}{\eq[t_1][t_2]}{$\Sat{M}{\indfrm}[s]$ तेव्हा आणि
  केवळ तेव्हाच, जेव्हा $\Value{t_1}{M}[s] = \Value{t_2}{M}[s]$.}

\tagitem{prvNot}{%
  \indcase{!A}{\lnot !B}{$\Sat{M}{\indfrm}[s]$ तेव्हा आणि केवळ तेव्हाच,
    जेव्हा $\Sat/{M}{!B}[s]$.}}{}

\tagitem{prvAnd}{%
  \indcase{!A}{(!B \land !C)}{$\Sat{M}{\indfrm}[s]$ तेव्हा आणि केवळ
    तेव्हाच, जेव्हा $\Sat{M}{!B}[s]$ आणि $\Sat{M}{!C}[s]$.}}{}

\tagitem{prvOr}{%
  \indcase{!A}{(!B \lor !C)}{$\Sat{M}{\indfrm}[s]$ तेव्हा आणि केवळ
    तेव्हाच, जेव्हा $\Sat{M}{!B}[s]$ किंवा $\Sat{M}{!C}[s]$ (किंवा
    दोन्ही).}}{}

\tagitem{prvIf}{%
  \indcase{!A}{(!B \lif !C)}{$\Sat{M}{\indfrm}[s]$ तेव्हा आणि केवळ
    तेव्हाच, जेव्हा $\Sat/{M}{!B}[s]$ किंवा $\Sat{M}{!C}[s]$ (किंवा
    दोन्ही).}}{}

\tagitem{prvIff}{%
  \indcase{!A}{(!B \liff !C)}{$\Sat{M}{\indfrm}[s]$ तेव्हा आणि केवळ
    तेव्हाच, जेव्हा एकतर $\Sat{M}{!B}[s]$ आणि $\Sat{M}{!C}[s]$ दोन्ही,
    किंवा $\Sat{M}{!B}[s]$ व $\Sat{M}{!C}[s]$ यांपैकी एकही नाही.}}{}

\tagitem{prvAll}{%
  \indcase{!A}{\lforall[x][!B]}{$\Sat{M}{\indfrm}[s]$ तेव्हा आणि केवळ
    तेव्हाच, जेव्हा प्रत्येक !!{element}~$m \in \Domain M$ साठी
    $\Sat{M}{!B}[\Subst{s}{m}{x}]$.}}{}

\tagitem{prvEx}{%
  \indcase{!A}{\lexists[x][!B]}{$\Sat{M}{\indfrm}[s]$ तेव्हा आणि केवळ
  तेव्हाच, जेव्हा किमान एका !!{element}~$m \in \Domain M$ साठी
  $\Sat{M}{!B}[\Subst{s}{m}{x}]$.}}{}
\end{enumerate}
\end{defn}

\begin{explain}
शेवटच्या \iftag{notprvEx,notprvAll}{उपवाक्यात}{दोन उपवाक्यांत}
चर-मूल्यांकने महत्त्वाची आहेत.\iftag{prvAll}{ ``प्रत्येक $m \in
\Domain{M}$ साठी $\Sat{M}{!B(m)}$'' असे म्हणून आपण
$\lforall[x][!B(x)]$ ची पूर्ति परिभाषित करू शकत नाही.}{}
\iftag{prvEx}{ ``किमान एका $m \in \Domain{M}$ साठी $\Sat{M}{!B(m)}$''
असे म्हणून आपण $\lexists[x][!B(x)]$ ची पूर्ति परिभाषित करू शकत नाही.}{}
कारण $m \in \Domain M$ असेल, तर ते भाषेचे चिन्ह नसते आणि म्हणून
$!B(m)$ हे !!a{formula} नसते (म्हणजे $\Subst{!B}{m}{x}$ अपरिभाषित
असते).~$\Struct{M}$ मधील प्रत्येक !!{element} चे नाव देणारी
!!{constant}s किंवा पदे
आपल्याजवळ उपलब्ध आहेत असेही गृहीत धरता येत नाही, कारण !!{structure}s
च्या व्याख्येत तशी कोणतीही अट नाही. प्रमाण भाषेतील !!{constant}s चा
संच !!{denumerable} आहे; म्हणून $\Domain{M}$ हा !!{enumerable} नसेल,
तर प्रत्येक वस्तूला नाव देण्यासाठी पुरेशी !!{constant}s सुद्धा नसतात.

ही अडचण सोडवण्यासाठी आपण !!{variable} ची मूल्यांकने आणतो; त्यांमुळे
चरे प्रांतातील !!{element}s शी थेट जोडता येतात. मग, उदाहरणार्थ,
\iftag{prvEx}{$\lexists[x][!B(x)]$}{$\lforall[x][!B(x)]$} ची~$\Struct M$
मध्ये पूर्ति तेव्हा आणि केवळ तेव्हाच होते, जेव्हा
\iftag{prvEx}{किमान एका $m \in \Domain{M}$ साठी $\Sat{M}{!B(m)}$}{प्रत्येक
$m \in \Domain{M}$ साठी $\Sat{M}{!B(m)}$}, असे न म्हणता आपण म्हणतो:
तिची~$\Struct M$ मध्ये~$s$ च्या \emph{सापेक्ष} पूर्ति तेव्हा आणि केवळ
तेव्हाच होते, जेव्हा $!B(x)$ ची~$\Subst{s}{m}{x}$ च्या सापेक्ष
\iftag{prvEx}{किमान एका}{प्रत्येक} $m \in \Domain M$ साठी पूर्ति होते.
%READERNOTE{OLFOL-013: गोठवलेल्या इंग्रजी मजकुराच्या अस्तित्ववाची शाखेत पूर्तिची आवश्यक अट वगळली आहे. समांतर वैश्विक शाखा आणि औपचारिक व्याख्या यांच्याशी जुळणारी अट मराठीत पुनःस्थापित केली आहे.}
\end{explain}

\begin{ex}
$\Lang{L} = \{a, b, f, R\}$ असू द्या, जिथे $a$ आणि $b$ ही
!!{constant}s, $f$ हे दोन-स्थानी !!{function} आणि $R$ हे दोन-स्थानी
!!{predicate} आहे. पुढीलप्रमाणे परिभाषित केलेली !!{structure}~$\Struct{M}$
विचारात घ्या:
\begin{enumerate}
\item $\Domain M = \{1, 2, 3, 4\}$
\item $\Assign{a}{M} = 1$
\item $\Assign{b}{M} = 2$
\item $x+y \le 3$ असेल तर $\Assign{f}{M}(x, y) = x+y$, अन्यथा $= 3$.
\item $\Assign{R}{M} = \{\tuple{1, 1}, \tuple{1, 2}, \tuple{2, 3}, \tuple{2, 4}\}$
\end{enumerate}
प्रत्येक !!{variable} ला $1 \in \Domain{M}$ नेमून देणारे फलन
$s(x) = 1$ हे~$\Struct{M}$ साठीचे चर-मूल्यांकन आहे.

मग
\begin{align*}
\Value{f(a,b)}{M}[s] & = \Assign{f}{M}(\Value{a}{M}[s], \Value{b}{M}[s]).
\intertext{$a$ आणि $b$ ही !!{constant}s असल्यामुळे $\Value{a}{M}[s]
  = \Assign{a}{M} = 1$ आणि $\Value{b}{M}[s] = \Assign{b}{M} = 2$. म्हणून}
\Value{f(a,b)}{M}[s] & = \Assign{f}{M}(1, 2) = 1+2 = 3.
\intertext{$f(f(a,b),a)$ चे मूल्य काढण्यासाठी पुढील अभिव्यक्ती विचारात घ्यावी लागते:}
\Value{f(f(a,b),a)}{M}[s] & = \Assign{f}{M}(\Value{f(a, b)}{M}[s],
\Value{a}{M}[s]) = \Assign{f}{M}(3, 1) = 3,
\intertext{कारण $3+1 > 3$. $s(x) = 1$ आणि $\Value{x}{M}[s] =
  s(x)$ असल्यामुळे आपल्याकडे पुढील समीकरणही आहे:}
\Value{f(f(a,b),x)}{M}[s] & = \Assign{f}{M}(\Value{f(a, b)}{M}[s],
\Value{x}{M}[s]) = \Assign{f}{M}(3, 1) = 3,
\end{align*}

आण्विक !!{formula}~$R(t_1, t_2)$ च्या फलसाधकांच्या मूल्यांची क्रमित
जोडी, म्हणजे $\tuple{\Value{t_1}{M}[s], \Value{t_2}{M}[s]}$, ही
$\Assign{R}{M}$ ची !!a{element} असेल, तर त्या सूत्राची पूर्ति होते.
म्हणून, उदाहरणार्थ, $\tuple{\Value{b}{M}, \Value{f(a,b)}{M}} =
\tuple{2, 3} \in \Assign{R}{M}$ असल्यामुळे
$\Sat{M}{R(b,f(a,b))}[s]$; पण $\tuple{1, 3} \notin \Assign{R}{M}$
असल्यामुळे $\Sat/{M}{R(x, f(a,b))}[s]$.
%READERNOTE{OLFOL-012: गोठवलेल्या इंग्रजी मजकुरात संबंधाचे अर्थनिर्धारण चर-मूल्यांकनावर अवलंबून असल्यासारखा अनावश्यक निर्देश जोडला आहे. रचनेतील संबंध स्थिर असल्यामुळे तो निर्देश मराठीत काढला आहे.}

अनाण्विक सूत्र~$!A$ ची पूर्ति होते का हे ठरवण्यासाठी त्याच्या मुख्य
संयोजकाला लागू होणारे विगमनाधारित व्याख्येतील उपवाक्य वापरावे. उदाहरणार्थ,
$R(a, a) \lif (R(b, x) \lor R(x, b))$ मधील मुख्य संयोजक~$\lif$ आहे आणि
\begin{align*}
  & \Sat{M}{R(a, a) \lif (R(b, x) \lor R(x, b))}[s] \text{ तेव्हा आणि केवळ तेव्हाच, जेव्हा }\\
  & \qquad
  \Sat/{M}{R(a,a)}[s] \text{ किंवा } \Sat{M}{R(b, x) \lor R(x, b)}[s]
  \intertext{$\Sat{M}{R(a,a)}[s]$ असल्यामुळे (कारण $\tuple{1,1} \in
    \Assign{R}{M}$), उत्तर अजून ठरवता येत नाही; प्रथम
    $\Sat{M}{R(b, x) \lor R(x, b)}[s]$ आहे का ते शोधावे लागते:}
  & \Sat{M}{R(b, x) \lor R(x, b)}[s] \text{ तेव्हा आणि केवळ तेव्हाच, जेव्हा }\\
  & \qquad \Sat{M}{R(b,x)}[s] \text{ किंवा } \Sat{M}{R(x, b)}[s]
  \intertext{आणि हे खरे आहे, कारण $\Sat{M}{R(x, b)}[s]$
    (कारण $\tuple{1,2} \in \Assign{R}{M}$).}
\end{align*}

आठवा, $s$ चा $x$-पर्याय म्हणजे~$s$ पासून फार तर~$x$ ला नेमून
दिलेल्या मूल्यातच भिन्न असणारे चर-मूल्यांकन. $\Domain{M}$ च्या प्रत्येक
!!{element} साठी~$s$ चा एक $x$-पर्याय आहे:
\begin{align*}
  s_1 & = \Subst{s}{1}{x}, &
  s_2 & = \Subst{s}{2}{x},\\
  s_3 & = \Subst{s}{3}{x}, &
  s_4 & = \Subst{s}{4}{x}.
\end{align*}
म्हणून, उदाहरणार्थ, $s_2(x) = 2$ आणि~$x$ शिवाय इतर सर्व चर~$y$ साठी
$s_2(y) = s(y) = 1$. $\Domain{M} = \{1, 2, 3, 4\}$ असल्यामुळे हेच
$\Struct{M}$ या रचनेसाठीचे~$s$ चे सर्व $x$-पर्याय आहेत. विशेषतः,
$s_1 = s$ हे लक्षात घ्या ($s$ हा नेहमीच स्वतःचा $x$-पर्याय असतो).

\iftag{prvEx}{अस्तित्ववाची संख्यापक असलेल्या
  !!{formula}~$\lexists[x][!A(x)]$ ची पूर्ति होते का हे ठरवण्यासाठी,
  किमान एका $m \in \Domain M$ साठी
  $\Sat{M}{!A(x)}[\Subst{s}{m}{x}]$ आहे का हे ठरवावे लागते. म्हणून,
  \[
  \Sat{M}{\lexists[x][(R(b,x) \lor R(x,b))]}[s],
  \]
  कारण $\Sat{M}{R(b,x) \lor R(x, b)}[\Subst{s}{1}{x}]$
  ($\Subst{s}{3}{x}$ हेसुद्धा चालले असते). पण,
  \[
  \Sat/{M}{\lexists[x][(R(b,x) \land R(x,b))]}[s]
  \]
  कारण कोणताही $m \in \Domain{M}$ निवडला, तरी
  $\Sat/{M}{R(b,x) \land R(x,b)}[\Subst{s}{m}{x}]$.}{}

\iftag{prvAll}{वैश्विक संख्यापक असलेल्या
  !!{formula}~$\lforall[x][!A(x)]$ ची पूर्ति होते का हे ठरवण्यासाठी,
  प्रत्येक $m \in \Domain M$ साठी
  $\Sat{M}{!A(x)}[\Subst{s}{m}{x}]$ आहे का हे ठरवावे लागते. म्हणून,
  \[
  \Sat{M}{\lforall[x][(R(x,a) \lif R(a,x))]}[s],
  \]
  कारण प्रत्येक $m \in \Domain M$ साठी
  $\Sat{M}{R(x,a) \lif R(a,x)}[\Subst{s}{m}{x}]$. $m = 1$ साठी
  $\Sat{M}{R(a,x)}[\Subst{s}{1}{x}]$ असल्यामुळे उत्तरवर्ती सत्य आहे;
  $m = 2$, $3$ आणि~$4$ साठी $\Sat/{M}{R(x,a)}[\Subst{s}{m}{x}]$
  असल्यामुळे पूर्ववर्ती असत्य आहे. पण,
  \[
  \Sat/{M}{\lforall[x][(R(a,x) \lif R(x,a))]}[s]
  \]
  कारण $\Sat/{M}{R(a,x) \lif R(x,a)}[\Subst{s}{2}{x}]$ (कारण
  $\Sat{M}{R(a, x)}[\Subst{s}{2}{x}]$ आणि $\Sat/{M}{R(x,
  a)}[\Subst{s}{2}{x}]$).}{}

\iftag{defEx}{अस्तित्ववाची संख्यापक असलेल्या
  !!{formula}~$\lexists[x][!A(x)]$ ची पूर्ति होते का हे ठरवण्यासाठी,
  $\Sat{M}{\lnot \lforall[x][\lnot !A(x)]}[s]$ आहे का हे ठरवावे लागते.
  उदाहरणार्थ, आपल्याकडे
  \[
  \Sat{M}{\lexists[x][(R(b,x) \lor R(x,b))]}[s].
  \]
  प्रथम, $\Sat{M}{R(b,x) \lor R(x, b)}[\Subst{s}{1}{x}]$
  ($\Subst{s}{3}{x}$ हेसुद्धा चालले असते). म्हणून,
  $\Sat/{M}{\lnot(R(b,x) \lor R(x,b))}[\Subst{s}{1}{x}]$ आणि त्यामुळे
  $\Sat/{M}{\lforall[x][\lnot(R(b,x) \lor R(x,b))]}[s]$; म्हणूनच
  $\Sat{M}{\lnot\lforall[x][\lnot(R(b,x) \lor R(x,b))]}[s]$. दुसरीकडे,
  \[
  \Sat/{M}{\lexists[x][(R(b,x) \land R(x,b))]}[s].
  \]
  याचे कारण $\Sat{M}{\lforall[x][\lnot(R(b,x) \land
  R(x,b))]}[s]$: कोणत्याही $m \in \Domain M$ साठी
  $\Sat{M}{R(b,x) \land R(x,b)}[\Subst{s}{m}{x}]$ नाही. या उदाहरणांवरून
  कदाचित तुमच्या लक्षात आले असेल की, किमान एका $m \in \Domain M$ साठी
  $\Sat{M}{!A(x)}[\Subst{s}{m}{x}]$ असेल तेव्हा आणि केवळ तेव्हाच
  $\Sat{M}{\lexists[x][!A(x)]}[s]$.
  %READERNOTE{OLFOL-014: गोठवलेल्या इंग्रजी उदाहरणात नकारानंतर दोन ठिकाणी एक जादा उघडा कंस आणि एका सूत्रात जादा स्वल्पविराम आहे. समांतर सूत्रे व व्याख्या यांनुसार मराठी आवृत्तीत ते दुरुस्त केले आहेत.}
}{}

\iftag{defAll}{वैश्विक संख्यापक असलेल्या
  !!{formula}~$\lforall[x][!A(x)]$ ची पूर्ति होते का हे ठरवण्यासाठी,
  $\Sat{M}{\lnot\lexists[x][\lnot !A(x)]}[s]$ आहे का हे ठरवावे लागते.
  उदाहरणार्थ,
  \[
  \Sat{M}{\lforall[x][(R(x,a) \lif R(a,x))]}[s],
  \]
  प्रथम, प्रत्येक $m \in \Domain M$ साठी
  $\Sat{M}{R(x,a) \lif R(a,x)}[\Subst{s}{m}{x}]$
  ($\Sat{M}{R(a,x)}[\Subst{s}{1}{x}]$ आणि $m = 2$, $3$ किंवा~$4$ साठी
  $\Sat/{M}{R(x,a)}[\Subst{s}{m}{x}]$). म्हणून असा कोणताही
  $m \in \Domain M$ नाही की
  $\Sat{M}{\lnot(R(x,a) \lif R(a,x))}[\Subst{s}{m}{x}]$; परिणामी
  $\Sat/{M}{\lexists[x][\lnot(R(x,a) \lif R(a,x))]}[s]$. म्हणून,
  $\Sat{M}{\lnot\lexists[x][\lnot(R(x,a) \lif R(a,x))]}[s]$.
  दुसरीकडे,
  \[
  \Sat/{M}{\lforall[x][(R(a,x) \lif R(x,a))]}[s],
  \]
  कारण $\Sat/{M}{R(a,x) \lif R(x,a)}[\Subst{s}{2}{x}]$, आणि म्हणून
  $\Sat{M}{\lexists[x][\lnot(R(a,x) \lif R(x,a))]}[s]$. या उदाहरणांवरून
  कदाचित तुमच्या लक्षात आले असेल की, प्रत्येक $m \in \Domain M$ साठी
  $\Sat{M}{!A(x)}[\Subst{s}{m}{x}]$ असेल तेव्हा आणि केवळ तेव्हाच
  $\Sat{M}{\lforall[x][!A(x)]}[s]$.
  %READERNOTE{OLFOL-015: गोठवलेल्या इंग्रजी उदाहरणात m = 2, 3, 4 साठी उत्तरवर्ती R(a,x) असत्य म्हटले आहे; m = 2 साठी ते सत्य आहे. सशर्त सूत्राचा पूर्ववर्ती R(x,a) या तिन्ही मूल्यांसाठी असत्य असल्यामुळे मराठीत तोच वापरला आहे.}
}{}

अधिक जटिल उदाहरणासाठी पुढील सूत्र विचारात घ्या:
\[
\lforall[x][(R(a,x) \lif \lexists[y][R(x,y)])].
\]
$\Sat/{M}{R(a,x)}[\Subst{s}{3}{x}]$ आणि
$\Sat/{M}{R(a,x)}[\Subst{s}{4}{x}]$ असल्यामुळे, सशर्त सूत्राच्या
उत्तरवर्तीची काळजी घ्यावी लागणारी रंजक प्रकरणे फक्त $m = 1$ आणि
$m = 2$ ही आहेत. $\Sat{M}{\lexists[y][R(x,y)]}[\Subst{s}{1}{x}]$ आहे
का? किमान एक $n \in \Domain M$ असा असेल की
$\Sat{M}{R(x,y)}[\Subst{\Subst{s}{1}{x}}{n}{y}]$, तर ते खरे आहे.
वस्तुतः, $n = 1$ घेतल्यास $\Subst{\Subst{s}{1}{x}}{n}{y} =
\Subst{s}{1}{y} = s$. $s(x) = 1$, $s(y) = 1$ आणि
$\tuple{1,1} \in \Assign{R}{M}$ असल्यामुळे उत्तर होकारार्थी आहे.
%READERNOTE{OLFOL-016: गोठवलेल्या इंग्रजी मजकुरात दुसऱ्या रंजक प्रकरणातील चलाचे m हे नाव गाळले आहे. संदर्भ आणि लगेच पुढील गणना यांनुसार ते मराठीत पुनःस्थापित केले आहे.}

$\Sat{M}{\lexists[y][R(x,y)]}[\Subst{s}{2}{x}]$ आहे का हे ठरवण्यासाठी
आपल्याला !!{variable} ची मूल्यांकने
$\Subst{\Subst{s}{2}{x}}{n}{y}$ तपासावी लागतात. इथे $n = 1$ साठी हे
मूल्यांकन $s_2 = \Subst{s}{2}{x}$ आहे आणि ते $R(x,y)$ ची पूर्ति करत
नाही ($s_2(x) = 2$, $s_2(y) = 1$ आणि
$\tuple{2,1}\notin \Assign{R}{M}$). पण
$\Subst{\Subst{s}{2}{x}}{3}{y} = \Subst{s_2}{3}{y}$ विचारात घ्या.
$\tuple{2,3} \in \Assign{R}{M}$ असल्यामुळे
$\Sat{M}{R(x,y)}[\Subst{s_2}{3}{y}]$ आणि म्हणून
$\Sat{M}{\lexists[y][R(x,y)]}[s_2]$.

म्हणून, प्रत्येक $m \in \Domain M$ साठी एकतर
$\Sat/{M}{R(a,x)}[\Subst{s}{m}{x}]$ ($m = 3$, $4$ असेल तर) किंवा
$\Sat{M}{\lexists[y][R(x,y)]}[\Subst{s}{m}{x}]$ ($m = 1$, $2$ असेल तर),
आणि म्हणून
\[
\Sat{M}{\lforall[x][(R(a,x) \lif \lexists[y][R(x,y)])]}[s].
\]
दुसरीकडे,
\[
\Sat/{M}{\lexists[x][(R(a,x) \land \lforall[y][R(x,y)])]}[s].
\]
फक्त $m = 1$ आणि $m = 2$ यांच्यासाठी
$\Sat{M}{R(a,x)}[\Subst{s}{m}{x}]$. पण~$m$ च्या या दोन्ही मूल्यांपैकी
प्रत्येकासाठी एक $n \in \Domain M$ आहे---पहिल्या मूल्यासाठी $n = 4$
आणि दुसऱ्यासाठी $n = 1$---ज्यामुळे
$\Sat/{M}{R(x,y)}[\Subst{\Subst{s}{m}{x}}{n}{y}]$ आणि म्हणून $m = 1$
व $m = 2$ साठी
$\Sat/{M}{\lforall[y][R(x,y)]}[\Subst{s}{m}{x}]$. एकंदरीत, असा कोणताही
$m \in \Domain M$ नाही की
$\Sat{M}{R(a,x) \land \lforall[y][R(x,y)]}[\Subst{s}{m}{x}]$.
%READERNOTE{OLFOL-017: गोठवलेल्या इंग्रजी सारांशात बाह्य वैश्विक चलाला चुकून n म्हटले आहे. सर्व चार बाह्य मूल्यांचा संदर्भ m नेच असल्यामुळे मराठीत m वापरले आहे.}
%READERNOTE{OLFOL-018: गोठवलेल्या इंग्रजी मजकुरात m = 1 आणि m = 2 दोन्हीसाठी n = 4 हा प्रतिदृष्टांत दिला आहे; पण (2,4) हा R मध्ये आहे. मराठीत अनुक्रमे n = 4 आणि n = 1 हे योग्य प्रतिदृष्टांत दिले आहेत.}
\end{ex}

\iftag{defEx}{%
\begin{prop}\ollabel{prop:sat-ex}
$\Sat{M}{\lexists[x][!B(x)]}[s]$ तेव्हा आणि केवळ तेव्हाच, जेव्हा~$s$ चा
एक $x$-पर्याय $s'$ असा आहे की $\Sat{M}{!B(x)}[s']$.
\end{prop}

\begin{proof}
  सराव.
\end{proof}
}{}

\tagprob{defEx}
\begin{prob}
  \olref[fol][syn][sat]{prop:sat-ex} सिद्ध करा.
\end{prob}
\tagendprob

\iftag{defAll}{%
\begin{prop}\ollabel{prop:sat-all}
$\Sat{M}{\lforall[x][!B(x)]}[s]$ तेव्हा आणि केवळ तेव्हाच, जेव्हा~$s$ च्या
प्रत्येक $x$-पर्याय~$s'$ साठी $\Sat{M}{!B(x)}[s']$.
\end{prop}

\begin{proof}
  सराव.
\end{proof}
}{}

\tagprob{defAll}
\begin{prob}
  \olref[fol][syn][sat]{prop:sat-all} सिद्ध करा.
\end{prob}
\tagendprob

\begin{prob}
एक !!{constant}, एक एक-स्थानी !!{function} आणि एक दोन-स्थानी
!!{predicate} असलेली $\Lang L = \{c, f, A\}$ असू द्या; आणि
!!{structure}~$\Struct{M}$ पुढीलप्रमाणे दिलेली असू द्या:
\begin{enumerate}
\item $\Domain M = \{1, 2, 3\}$
\item $\Assign{c}{M} = 3$
\item $\Assign{f}{M}(1) = 2, \Assign{f}{M}(2) = 3, \Assign{f}{M}(3) = 2$
\item $\Assign{A}{M} = \{\tuple{1, 2}, \tuple{2, 3}, \tuple{3, 3}\}$
\end{enumerate}
(a) सर्व !!{variable}s~$v$ साठी $s(v) = 1$ असू द्या. पुढील विधान
खरे आहे की नाही हे शोधा:
\[
\Sat{M}{\lexists[x][(A(f(z), c) \lif \lforall[y][(A(y, x) \lor A(f(y),
      x))])]}[s]
\]
तुमचे उत्तर का योग्य आहे ते स्पष्ट करा.

(b) ज्यात या !!{formula} ची पूर्ति होत नाही अशी वेगळी रचना आणि
!!{variable} चे मूल्यांकन द्या.
\end{prob}

\end{document}
